%==============================================================================
% ADDITIONAL TABLES
%
% Three tables that a reviewer is likely to ask for, in descending order of
% necessity. Each defends a decision the paper makes rather than adding
% description for its own sake.
%
%   Table C.1  branch scope      -> §2 Methods  (NECESSARY: defends a choice)
%   Table C.2  selection funnel  -> §2.1        (makes selection reproducible)
%   Table C.3  review actors     -> §3.3        (the split asked for explicitly)
%
% Requires: booktabs. Table C.1 also uses \multicolumn and \cmidrule.
%==============================================================================


%==============================================================================
% TABLE C.1 — BRANCH SCOPE.  Put this in §2, near the branch-scope paragraph.
% This is the one I would not omit: it is the evidence for analysing a branch
% other than the repository default, which a reviewer will question.
%==============================================================================
\begin{table}[htbp]
\centering
\caption{Commits reachable from each branch, January 2025 to August 2026.
Counts overlap, since a commit merged from \texttt{develop} into
\texttt{master} is reachable from both. The default branch is \texttt{master}
for 19 of the 20 repositories and \texttt{main} for one.}
\label{tab:branches}
\small
\begin{tabular}{@{}lrrr@{}}
\toprule
Repository & \texttt{master} & \texttt{develop} & all branches \\
\midrule
esignet                     & 461 & 539 & 1165 \\
id-repository               & 131 & 294 &  808 \\
partner-management-services &  64 & 442 &  731 \\
esignet-signup              & 303 & 323 &  544 \\
admin-services              & 179 & 141 &  422 \\
registration                &  69 &  61 &  390 \\
resident-services           & 107 & 310 &  383 \\
commons                     &  94 & 129 &  383 \\
id-authentication           &  71 & 146 &  370 \\
keymanager                  &  71 &  70 &  251 \\
pre-registration            &  73 &  63 &  231 \\
android-registration-client & 100 & 155 &  218 \\
bio-utils                   &  87 &  59 &  190 \\
print                       &  63 &  89 &  176 \\
packet-manager              &  41 &  95 &  168 \\
mosip-openid-bridge         &  59 &  33 &  156 \\
registration-client         &   2 &  39 &   79 \\
nfiq                        &  43 &  27 &   74 \\
imagedecoder                &   7 &  14 &   51 \\
gist                        &   0 &   0 &    1 \\
\midrule
\textbf{Total} & \textbf{2025} & \textbf{3029} & \textbf{6791} \\
\textit{share of all branches} & \textit{30\%} & \textit{45\%} & \textit{100\%} \\
\bottomrule
\end{tabular}
\end{table}

% --- accompanying prose, to replace the "Branch scope" paragraph in §2.5 -----
%
% \paragraph{Branch scope.} MOSIP follows gitflow, and only 30\% of commits in
% the window are reachable from the repository default branch
% (Table~\ref{tab:branches}). A further 15\% are reachable from
% \texttt{develop}, and the remaining 55\% sit on feature branches that had not
% been merged into either by the end of the window. No single branch therefore
% captures development activity completely, and the choice of branch is a
% substantive one rather than a formality.
%
% The two branches behave differently across repositories: \texttt{develop}
% carries far more activity in some (\texttt{partner-management-services}, 442
% against 64) while \texttt{master} carries more in others
% (\texttt{admin-services}, 179 against 141), reflecting differences in release
% discipline rather than a uniform convention.
%
% We report velocity on two scopes. The all-branches count answers "how much
% work was performed", and is the figure used in the headline descriptive
% statistics. The \texttt{develop} count answers "how much work reached the
% integration branch", and matches the tree the SonarQube snapshots analyse, so
% it is the scope used wherever velocity and quality are compared.
% Table~\ref{tab:its} reports the interrupted time series under both; no
% quality coefficient differs, and the velocity coefficients differ in
% magnitude but not in sign or significance.


%==============================================================================
% TABLE C.2 — SELECTION FUNNEL.  Put this in §2.1.
% Makes the selection reproducible and pre-empts "why these twenty".
%==============================================================================
\begin{table}[htbp]
\centering
\caption{Repository selection. Each row applies to the survivors of the row
above.}
\label{tab:selection}
\small
\begin{tabular}{@{}lrl@{}}
\toprule
Step & Remaining & Criterion \\
\midrule
Public repositories in the \texttt{mosip} organisation & 109 & --- \\
Not archived, not a fork                               &  97 & active and original \\
Primary language on the JVM                            &  48 & Java, Kotlin, Scala, Groovy \\
Not a mock, test harness, or compliance toolkit        &  41 & name-based exclusion \\
Not infrastructure, deployment or configuration        &  38 & name-based exclusion \\
Not documentation or a demonstration artefact          &  36 & name-based exclusion \\
\midrule
\textbf{Top 20 by JVM source volume}                   & \textbf{20} & \textbf{analysed} \\
\bottomrule
\end{tabular}
\end{table}


%==============================================================================
% TABLE C.3 — REVIEW ACTORS.  Put this in §3.3, before the coverage table.
% Addresses the request to separate AI from human from other automation.
%==============================================================================
\begin{table}[htbp]
\centering
\caption{Review submissions by actor, in the six months either side of
review-bot adoption. Non-AI automation (Dependabot, CI and release bots) opens
pull requests in MOSIP but almost never reviews them, so the human/AI split is
effectively exhaustive.}
\label{tab:review-actors}
\small
\begin{tabular}{@{}lrrrr@{}}
\toprule
 & \multicolumn{2}{c}{Pre ($\tau \in [-6,-1]$)} & \multicolumn{2}{c}{Post ($\tau \in [+1,+6]$)} \\
\cmidrule(lr){2-3}\cmidrule(lr){4-5}
Reviewer & submissions & share & submissions & share \\
\midrule
Human              & 2516 & 99.8\% & 3222 & 57.9\% \\
AI review bot      &    0 &  0.0\% & 2339 & 42.1\% \\
Other automation   &    4 &  0.2\% &    0 &  0.0\% \\
\midrule
\textbf{Total}     & \textbf{2520} & & \textbf{5561} & \\
\textit{Pull requests opened} & \textit{1731} & & \textit{1441} & \\
\bottomrule
\end{tabular}
\end{table}
